We created a comparison methodology based on various papers that compare LLMs.
This way, we came-up with a methodology that takes multiple aspects into consideration, not only a few.
The methodology details how every step should be performed, and what must be taken into consideration.

The methodology starts with the right prompt selections.
In contrast to methods we all got used to, LLMs cannot be used without proper prompting.
The first thing we must do is to select the right prompt for the model, for the task we are going to use the model for.

After selecting the right prompt, we must evaluate the functional validity of the models.
Functional validity, simply said, measures whether the model can provide code that does what it should.
The hardest task at this point is to find a way to measure that it does ''what it should do''.
This is most likely to be performed via automated tests, although it is important to note that these tests might not take the generated code as is, so a framework might be needed.
The use of frameworks depends entirely on the requirements because they may require that the code be testable in existing tests.

Getting the tests is not an easy task either, as it is either up to the team to create the tests or they must use existing benchmarks; although in that case, overfitting must be taken into consideration.
Either way, tests are not easy to obtain in good quality.

At the third step of the methodology, we must evaluate the non-functional validity or technical quality.
This step provides that the introduced LLM maintains the source code quality or rather increases it.

As our last step, We included the human evaluation, in order to see if developers are willing to accept source code generated by LLMs.
This step, although difficult and expensive to perform, provides a better overview if an LLM should be introduced to the developer team or not.

As an additional step, We done a case study to provide an actual example, how to use our methodology.
We also note, that additional evaluations might be added, such as complexity or performance tests, depending on the needs of the development team.
